\hypertarget{metaclasses-descriptor-protocol-and-type-slots}{%
\chapter{Metaclasses, Descriptor Protocol, and type
Slots}\label{metaclasses-descriptor-protocol-and-type-slots}}

\hypertarget{the-metaclass-class-creation-pipeline}{%
\subsection{25.1 The Metaclass Class Creation
Pipeline}\label{the-metaclass-class-creation-pipeline}}

In Python, classes are objects themselves. A \textbf{metaclass} is the
class of a class; it defines how classes are constructed. The default
metaclass for all objects is \passthrough{\lstinline!type!}.

\hypertarget{metaclass-lookup-and-construction-path}{%
\subsubsection{1. Metaclass Lookup and Construction
Path}\label{metaclass-lookup-and-construction-path}}

When CPython executes a class definition block:

\begin{lstlisting}[language=Python]
class MyClass(metaclass=MyMeta):
    x = 1
\end{lstlisting}

The VM resolves class creation by running these sequential steps: 1.
\textbf{Isolate Namespace}: The interpreter executes the class body code
inside a temporary namespace dictionary. 2. \textbf{Resolve Metaclass}:
It checks the inheritance tree for explicit metaclasses. If none are
declared, it defaults to \passthrough{\lstinline!type!}. 3.
\textbf{Execute \passthrough{\lstinline!\_\_new\_\_!}}: It invokes
\passthrough{\lstinline!MyMeta.\_\_new\_\_(meta, name, bases, namespace)!}.
Under the hood, this calls \passthrough{\lstinline!type.\_\_new\_\_!}
which allocates the \passthrough{\lstinline!PyTypeObject!} struct on the
heap, setting up C-level slots. 4. \textbf{Execute
\passthrough{\lstinline!\_\_init\_\_!}}: It initializes class state
parameters before returning the type object.

\begin{center}\rule{0.5\linewidth}{0.5pt}\end{center}

\hypertarget{the-descriptor-protocol-and-attribute-lookup-precedence}{%
\subsection{25.2 The Descriptor Protocol and Attribute Lookup
Precedence}\label{the-descriptor-protocol-and-attribute-lookup-precedence}}

A descriptor is an object that customizes attribute access behavior by
implementing methods in the descriptor protocol: *
\passthrough{\lstinline!\_\_get\_\_(self, instance, owner)!} *
\passthrough{\lstinline!\_\_set\_\_(self, instance, value)!} *
\passthrough{\lstinline!\_\_delete\_\_(self, instance)!}

\hypertarget{data-descriptors-vs.-non-data-descriptors}{%
\subsubsection{1. Data Descriptors vs.~Non-Data
Descriptors}\label{data-descriptors-vs.-non-data-descriptors}}

\begin{itemize}
\tightlist
\item
  \textbf{Data Descriptor}: Implements both
  \passthrough{\lstinline!\_\_get\_\_!} and
  \passthrough{\lstinline!\_\_set\_\_!} (or
  \passthrough{\lstinline!\_\_delete\_\_!}).
\item
  \textbf{Non-Data Descriptor}: Only implements
  \passthrough{\lstinline!\_\_get\_\_!} (typically used for methods).
\end{itemize}

\hypertarget{attribute-lookup-precedence-hierarchy}{%
\subsubsection{2. Attribute Lookup Precedence
Hierarchy}\label{attribute-lookup-precedence-hierarchy}}

When retrieving an attribute \passthrough{\lstinline!obj.name!}, CPython
resolves the lookup path in this strict order:

\begin{lstlisting}
                          [Lookup obj.name]
                                 |
           Does "name" exist in class MRO as a Data Descriptor?
                               /   \
                             Yes    No
                             /       \
         [Call descriptor __get__]  Does "name" exist in instance __dict__?
                                       /   \
                                     Yes    No
                                     /       \
                       [Return from __dict__]  Does "name" exist in class MRO
                                               as a Non-Data Descriptor?
                                                 /   \
                                               Yes    No
                                               /       \
                                   [Call descriptor]  Is there a class attribute?
                                                        /   \
                                                      Yes    No
                                                      /       \
                                         [Return value]   Raise AttributeError
\end{lstlisting}

\begin{center}\rule{0.5\linewidth}{0.5pt}\end{center}

\hypertarget{cpython-type-slots}{%
\subsection{25.3 CPython Type Slots}\label{cpython-type-slots}}

To optimize attribute access and function calls at the C level, CPython
uses \textbf{Type Slots}. Instead of looking up methods like
\passthrough{\lstinline!\_\_repr\_\_!} or
\passthrough{\lstinline!\_\_add\_\_!} in the class dictionary at
runtime, CPython populates static function pointers directly inside the
\passthrough{\lstinline!PyTypeObject!} struct definition:

\begin{lstlisting}[language=C]
typedef struct _typeobject {
    PyObject_VAR_HEAD
    const char *tp_name;                 /* For printing, in format <module>.<name> */
    Py_ssize_t tp_basicsize, tp_itemsize; /* For allocation */
    
    /* Type Slots (Static Function Pointers) */
    destructor tp_dealloc;               /* Deallocation handler */
    reprfunc tp_repr;                    /* Representation builder */
    getattrfunc tp_getattr;              /* Attribute lookup hook */
    setattrfunc tp_setattr;              /* Attribute assignment hook */
    
    /* Protocol Table Slots */
    PyNumberMethods *tp_as_number;       /* Math function pointers (e.g. nb_add) */
    PySequenceMethods *tp_as_sequence;   /* Indexing function pointers (e.g. sq_item) */
    PyMappingMethods *tp_as_mapping;     /* Map lookup function pointers (e.g. mp_subscript) */
} PyTypeObject;
\end{lstlisting}

When a descriptor or custom method is compiled, CPython populates these
function pointers. This allows the VM to execute operations (like
integer addition \passthrough{\lstinline!a + b!}) using direct C
function calls (\passthrough{\lstinline!tp\_as\_number->nb\_add(a, b)!})
instead of slow dictionary lookups.

\begin{center}\rule{0.5\linewidth}{0.5pt}\end{center}
